% !TeX spellcheck = en_US
\documentclass[french]{yLectureNote}

\title{Électromagnétisme 3}
\subtitle{Physique}
\author{Paulhenry Saux}
\date{\today}
\yLanguage{Français}
%jesse.groenen@cemes.fr
\professor{M.Brut}
\usepackage{graphicx}%-pour mettre des images
\usepackage[utf8]{inputenc}%encodage
\usepackage{geometry}%pour modifier les tailles et mettre a4paper
%\usepackage{awesomebox}%pour les boites d'exercices, de pbq et de croquis d\'esactiv\'e pour les TP de PC
\usepackage{tikz}%pour deiffner + d\'ependance de chemfig
% \usepackage{tabularx}%pour dimensionner automatiquement les tableaux avec variable X
\usepackage{awesomebox}%Pour les boites info, danger et autres
\usepackage{menukeys}%Pour deiffner les touches de Calculatrice
\usepackage{fancyhdr}%pour les en-t\^ete personnalis\'ees
\usepackage{blindtext}%pour les liens
\usepackage{hyperref}%pour les liens (\`a mettre en dernier)
\usepackage{caption}%pour la francisation de la l\'egende table vers Tableau
\usepackage{pifont}
\usepackage{array}%pour les tableaux
\usepackage{lipsum}
\usepackage{yFlatTable}
\usepackage{multicol}
\usepackage{tabularx}
\usepackage{booktabs}
\usepackage{esint}
\newcommand{\Lim}[1]{\lim\limits_{\substack{1}}\:}
\renewcommand{\vec}{\overrightarrow}
\newcommand{\N}[0]{\mathbb{N}}
\newcommand{\dd}{\mathrm{d}}
\newcommand{\norm}[1]{||\vec{1}||}
\newcommand{\fo}{\psi(\vec{r},t)}
\newcommand{\HH}{\hat{H}}
\newcommand{\hb}{\hbar}
\newcommand{\lap}{\nabla^2}
\newcommand{\lapcc}{\frac{\partial^2 }{\partial x^2}+\frac{\partial^2 }{\partial y^2}+\frac{\partial^2 }{\partial z^2}}
\newcommand{\E}{\vec{E}(M)}
\newcommand{\B}{\vec{B}(M)}
\newcommand{\V}{V(M)}
\newcommand{\er}{\vec{e_{\rho}}}
\newcommand{\ez}{\vec{e_{z}}}
\newcommand{\ep}{\vec{e_{\varphi}}}
\newcommand{\pip}{\pi^+}
\newcommand{\pim}{\pi^-}
\newcommand{\mv}{\mathcal{V}}
\DeclareMathOperator{\Div}{Div}
\newcommand{\rot}{\vec{\mathrm{rot}}}
\newcommand{\grad}{\vec{\mathrm{grad}}}
\newcommand{\vds}{\vec{\dd S}}
\newcommand{\Ea}{\vec{E_{aux}}}
%\setcounter{chapter}{3}
\begin{document}

\chapter{Propagation d'onde dans la matière}

\section{Notions de cours}

\begin{definition}[Permittivité et fonction diélectrique]
La fonction diélectrique relative correspond à la permittivité relative du milieu, notée $\underline{\epsilon_r}$. En l'absence d'absorption, cette grandeur est réelle : $\underline{\epsilon_r} = \epsilon_r$.
\end{definition}

\checkInfo{Régime statique}{Dans ce chapitre, la pulsation de l'onde vérifie $\omega \ll \omega_{\text{domaines d'absorption}}$. On se place donc dans l'approximation statique où la permittivité prend sa valeur statique, notée $\epsilon_s$, avec $\epsilon_s > 1$.}

\subsection{Équations de Maxwell dans un milieu localement neutre}

Les champs électrique $\vec{E}$ et magnétique $\vec{B}$ sont recherchés sous forme d'ondes monochromatiques de pulsation $\omega$. En notation complexe :
$$\underline{\vec{E}} = \underline{\vec{E_m}} e^{-i\omega t}$$

Dans un milieu localement neutre ($\rho_{\text{libre}} = 0$ et $\vec{j}_{\text{libre}} = \vec{0}$), les équations de Maxwell s'écrivent :
\begin{itemize}
    \item $\rot \underline{\vec{H}} = \frac{\partial \underline{\vec{D}}}{\partial t} = -i\omega \underline{\vec{D}}$
    \item $\rot \underline{\vec{E}} = -\frac{\partial \underline{\vec{B}}}{\partial t} = i\omega \underline{\vec{B}}$
    \item $\Div \underline{\vec{D}} = 0$
    \item $\Div \underline{\vec{B}} = 0$
\end{itemize}

Ces équations sont complétées par les deux relations constitutives du milieu non magnétique :
\begin{itemize}
    \item $\underline{\vec{D}} = \epsilon_0 \epsilon_s \underline{\vec{E}}$
    \item $\underline{\vec{H}} = \frac{\underline{\vec{B}}}{\mu_0}$
\end{itemize}

\subsection{Équation de propagation}

Pour établir l'équation de propagation de $\underline{\vec{E}}$, on prend le rotationnel de l'équation de Maxwell-Faraday en utilisant l'identité vectorielle : 
$$\rot(\rot \underline{\vec{E}}) = \grad(\Div \underline{\vec{E}}) - \lap \underline{\vec{E}}$$

Comme $\Div \underline{\vec{D}} = 0$ et que le milieu est homogène ($\epsilon_s = \text{cte}$), on a $\Div \underline{\vec{E}} = 0$, d'où $\rot(\rot \underline{\vec{E}}) = - \lap \underline{\vec{E}}$.

\begin{flalign*}
   \rot(\rot \underline{\vec{E}}) &= i\omega \rot\underline{\vec{B}} &\\
   &= i\omega \mu_0 \rot \underline{\vec{H}} &\\
   &= i\omega \mu_0 (-i\omega \underline{\vec{D}}) = \omega^2 \mu_0 \underline{\vec{D}} &\\
   &= \omega^2 \epsilon_0 \mu_0 \epsilon_s \underline{\vec{E}} &
\end{flalign*}

En utilisant la relation $c^2 = \frac{1}{\epsilon_0 \mu_0}$, on obtient l'équation de Helmholtz :
\begin{proposition}[Équation de Helmholtz]
$$\lap\underline{\vec{E}} + \frac{\omega^2}{c^2}\epsilon_s \underline{\vec{E}} = \vec{0}$$
\end{proposition}

\subsection{Solution et relation de dispersion}

On cherche une solution sous la forme d'une onde plane progressive harmonique (OPPH) se propageant selon l'axe $(Oz)$ :
$$\underline{\vec{E}} = \vec{E}_0 e^{i(\underline{k}z - \omega t)} \quad \text{avec} \quad \vec{E}_0 = E_0\vec{e_x}$$

\tipsInfo{Choix de la structure de l'onde}{L'onde est choisie transversale ($\vec{E} \cdot \vec{k} = 0$). Si $\vec{k} = \underline{k}\vec{e_z}$, le champ électrique vibre dans le plan $(Oxy)$. On choisit ici arbitrairement une polarisation rectiligne selon $\vec{e_x}$.}

En injectant cette forme dans l'équation de propagation, l'opérateur laplacien se réduit à $\lap \underline{\vec{E}} = -\underline{k}^2 \underline{\vec{E}}$. On obtient :
$$-\underline{k}^2 \underline{\vec{E}} + \frac{\omega^2}{c^2}\epsilon_s \underline{\vec{E}} = \vec{0}$$

\begin{theorem}[Relation de dispersion]
La relation de dispersion dans le milieu s'écrit :
$$\underline{k}^2 = \epsilon_s \frac{\omega^2}{c^2}$$
\end{theorem}

\subsubsection{Analyse de l'indice complexe}

On définit l'indice de réfraction complexe $\underline{n} = n + i\kappa$ tel que $\underline{n}^2 = \underline{\epsilon_r} = \epsilon_s$. En développant :
$$n^2 - \kappa^2 + 2in\kappa = \epsilon_s$$

Par identification des parties réelles et imaginaires :
$$n^2 - \kappa^2 = \epsilon_s \quad \text{et} \quad 2n\kappa = 0$$

Comme $\epsilon_s > 1$, la grandeur est strictement positive. La seule solution physique pour un milieu transparent est :
$$n = \sqrt{\epsilon_s} \quad \text{et} \quad \kappa = 0$$

Le vecteur d'onde est donc réel : $k = \frac{n\omega}{c}$. L'expression finale du champ est :
$$\underline{\vec{E}}(z,t) = E_0 e^{i\left(\frac{n\omega}{c}z - \omega t\right)}\vec{e_x}$$
L'amplitude réelle $E_0$ est constante au cours de la propagation : **il n'y a pas d'atténuation**.

\subsection{Propriétés de l'onde}

\begin{itemize}
    \item \textbf{Vitesse de phase :} La phase de l'onde est $\phi = kz - \omega t$. Par définition :
    $$v_{\varphi} = \frac{\omega}{k} = \frac{c}{n} = \frac{c}{\sqrt{\epsilon_s}} < c$$
    
    \item \textbf{Période spatiale (Longueur d'onde) :} 
    $$\lambda = \frac{2\pi}{k} = \frac{2\pi c}{n\omega} = \frac{\lambda_0}{n}$$
    où $\lambda_0 = \frac{2\pi c}{\omega}$ est la longueur d'onde dans le vide.
    
    \item \textbf{Champ magnétique de l'onde : } En utilisant la relation de structure des OPPH $\underline{\vec{B}} = \frac{\vec{k} \wedge \underline{\vec{E}}}{\omega}$ :
    $$\underline{\vec{B}} = \frac{k}{\omega}\vec{e_z} \wedge \left(\underline{E_x}\vec{e_x}\right) = \frac{n}{c}\underline{E_x}\vec{e_y}$$
\end{itemize}

\subsection{Vecteur de Poynting et Intensité}

Le vecteur de Poynting s'exprime à partir des champs réels : $\vec{\Pi} = \vec{E} \wedge \frac{\vec{B}}{\mu_0} = \frac{1}{\mu_0}E_x B_y \vec{e_z}$.

L'intensité $I$ correspond à la valeur moyenne temporelle du flux du vecteur de Poynting :
$$I = \langle \norm{\vec{\Pi}} \rangle = \frac{1}{2\mu_0} \text{Re}\left(\underline{E_x}\cdot \underline{B_y}^*\right) = \frac{1}{2\mu_0} E_0 B_0 = \frac{1}{2\mu_0} E_0^2 \frac{n}{c}$$

En utilisant $c^2 = \frac{1}{\epsilon_0 \mu_0} \implies \frac{1}{\mu_0 c} = \epsilon_0 c$, on obtient :
$$I = \frac{1}{2} n c \epsilon_0 E_0^2$$

 \textit{Application numérique à faire : Calculer $v_{\varphi}$, $\lambda$ et $I$ avec les données de l'énoncé.}

\section{L'arséniure de gallium (GaAs) dopé}

Dans un semi-conducteur dopé, on combine la polarisation du milieu pur (matrice cristalline) et celle apportée par les porteurs libres issus du dopage.

\subsection{Modèle de Drude pour les porteurs libres}

On applique la seconde loi de Newton à un électron de conduction de masse effective $m^*$ et de charge $-e$ :
$$m^* \frac{\dd \vec{v}}{\dd t} = \vec{F}_{\text{Lorentz}} - \frac{m^*}{\tau}\vec{v}$$

\checkInfo{Hypothèse du modèle}{On néglige ici les collisions, ce qui revient à poser un temps moyen entre deux chocs $\tau \to \infty$. Le terme de force de frottement fluide $\frac{m^*}{\tau}\vec{v}$ est donc négligé.}

\warningInfo{Électrons libres vs liés}{S'il s'agissait d'électrons liés (modèle de l'électron élastique de Lorentz), il aurait fallu rajouter une force de rappel élastique $-m^*\omega_0^2\vec{r}$. Ce terme est absent pour les porteurs libres du dopage.}

\subsection{Hiérarchie des forces}

Le rapport de la force magnétique sur la force électrique est donné par :
$$\frac{\norm{\vec{F}_m}}{\norm{\vec{F}_e}} = \frac{\norm{-e\vec{v}\wedge\vec{B}}}{\norm{-e\vec{E}}} \sim \frac{v B_0}{E_0} = \frac{v}{v_{\varphi}} = \frac{v}{c}n \ll 1$$
La force magnétique est donc totalement négligeable devant la force électrique en régime non relativiste.

\subsection{Équation du mouvement et conductivité}

L'équation du mouvement linéarisée en régime harmonique s'écrit :
$$-i\omega m^* \underline{\vec{v}} = -e \underline{\vec{E}} \implies \underline{\vec{v}} = \frac{-e}{i\omega m^*}\underline{\vec{E}}$$

Le vecteur densité de courant volumique associé aux $N$ porteurs libres par unité de volume vaut :
$$\underline{\vec{J}} = -e N \underline{\vec{v}} = -i \frac{N e^2}{m^* \omega}\underline{\vec{E}}$$

Par identification avec la loi d'Ohm locale $\underline{\vec{J}} = \underline{\sigma}\underline{\vec{E}}$, on en déduit la conductivité complexe :
$$\underline{\sigma}(\omega) = -i \frac{N e^2}{m^* \omega}$$

\subsection{Courant de polarisation et pulsation plasma}

Le déplacement des charges libres génère un moment dipolaire macroscopique, définissant une polarisation des charges libres $\underline{\vec{P}}_{\text{libre}}$ telle que $\underline{\vec{J}} = \frac{\partial \underline{\vec{P}}_{\text{libre}}}{\partial t} = -i\omega \underline{\vec{P}}_{\text{libre}}$. On a donc :
$$\underline{\vec{P}}_{\text{libre}} = \frac{i}{\omega}\underline{\vec{J}} = -\frac{N e^2}{m^* \omega^2}\underline{\vec{E}}$$

On introduit la pulsation plasma du milieu $\omega_p$, définie par : $\omega_p^2 = \frac{N e^2}{m^* \epsilon_0 \epsilon_s}$. On peut alors écrire :
$$\underline{\vec{P}}_{\text{libre}} = -\epsilon_0 \epsilon_s \frac{\omega_p^2}{\omega^2}\underline{\vec{E}}$$

\subsection{Polarisation totale et permittivité relative du milieu dopé}

La polarisation totale du milieu est la somme de la polarisation du cristal pur (GaAs non dopé, de permittivité $\epsilon_s$) et de celle des charges libres :
$$\underline{\vec{P}} = \underline{\vec{P}}_{\text{pur}} + \underline{\vec{P}}_{\text{libre}} = \epsilon_0(\epsilon_s - 1)\underline{\vec{E}} - \epsilon_0 \epsilon_s \frac{\omega_p^2}{\omega^2}\underline{\vec{E}}$$

$$\underline{\vec{P}} = \epsilon_0 \left[ \epsilon_s \left(1 - \frac{\omega_p^2}{\omega^2}\right) - 1 \right] \underline{\vec{E}} = \epsilon_0 (\underline{\epsilon_r} - 1)\underline{\vec{E}}$$

Par identification, la permittivité relative globale du milieu dopé vaut :
$$\epsilon_r(\omega) = \epsilon_s \left(1 - \frac{\omega_p^2}{\omega^2}\right)$$
Cette valeur est purement réelle : il n'y a pas de terme imaginaire, donc pas de dissipation d'énergie par effet Joule à cette approximation.

\subsection{Régimes de propagation et caractère de l'onde}

L'indice complexe vérifie $\underline{n}^2 = (n+i\kappa)^2 = \epsilon_r(\omega)$. Deux cas de figure se présentent selon le signe de $\epsilon_r$ (contrôlé par la pulsation $\omega$ vis-à-vis de la pulsation plasma $\omega_p$, et donc par le dopage $N$) :

\begin{table}[h!]
\centering
\begin{tabular}{|c|c|c|}
\hline
\textbf{Condition} & $N < N_0 \iff \omega > \omega_p$ & $N > N_0 \iff \omega < \omega_p$ \\ \hline
\textbf{Signe de $\epsilon_r$} & $\epsilon_r > 0$ & $\epsilon_r < 0$ \\ \hline
\textbf{Composantes de $\underline{n}$} & $n = \sqrt{\epsilon_r}$ et $\kappa = 0$ & $n = 0$ et $\kappa = \sqrt{|\epsilon_r|}$ \\ \hline
\textbf{Forme du champ $\underline{\vec{E}}$} & $E_0 e^{i(kz-\omega t)}\vec{e_x}$ & $E_0 e^{-\frac{\omega}{c}\kappa z} e^{-i\omega t}\vec{e_x}$ \\ \hline
\textbf{Nature de l'onde} & \textbf{Onde propagative} (Milieu transparent) & \textbf{Onde évanescente} (Milieu réfléchissant) \\ \hline
\end{tabular}
\caption{Régimes de propagation dans le GaAs dopé.}
\end{table}
\end{document}
